\begin{chaptersummary}
  \item Your service already performs the join federation performs. The
        difference is that three integers which lived inside one process now
        have to be serialized, which makes them fields, which makes them part
        of a contract two teams have to agree on. Everything in the model
        follows from that one crossing.
  \item A subgraph is an ordinary GraphQL service that also applies
        \code{@link} to its schema, defines the federation directives, and
        resolves \code{\_service} and \code{\_entities}. A supergraph is what a
        composer makes of several of them, and it is not a document anybody
        writes or owns. Only the first of those two words has a specification
        behind it.
  \item Composition and execution are different times with different failures.
        A field that has not worked since the deployment, for anybody, is
        composition. A field that works for some requests and not others is
        execution. They have separate tooling and separate people to ask.
  \item An entity is a type with a \code{@key}, and the key is a selection set
        written as a string. Declaring it forces every subgraph that touches
        the type to publish the key fields and to compute the same value for
        the same object, which nothing checks. \code{\_entities} takes those
        keys as untyped JSON and must answer in the order it was asked, because
        position is the only thing connecting an answer to its question.
  \item A field belongs to the subgraph that owns its data, and the composer
        enforces that as a claim rather than a convention: two subgraphs
        declaring one field fail composition unless both call it
        \code{@shareable}. Global object identification has already bought this
        book one of those, on \code{Query.node}.
\end{chaptersummary}
